% ===========================================================================
% The flagged cell is free of the noise-trading intensity.
% ===========================================================================

\providecommand{\Eop}{\mathbb{E}}
\providecommand{\Prb}{\mathbb{P}}
\providecommand{\ind}{\mathbf{1}}
\providecommand{\ND}{\mathrm{ND}}
\providecommand{\Sfl}{\mathsf{S}_F}

\begin{lemma}[The flagged cell does not move with liquidity]
\label{lem:flagged-kappa-free}
Fix the cutoff policy and the execution policy, and let the disclosure rule $(\tau,T)$ be fixed.
Assume the primitive draws are mutually independent with strictly positive variances; the menu, the
image of the coarsening $\Gamma$, the noise support and the calendar are finite; every stake path
$s\mapsto B_j(s,d)$ is Borel, with Voice stakes weakly increasing in the signal and in the date; the
blockholder does not revise the stake path in response to realised order flow or prices; the legal-clock definitions
hold, including $D=1\Rightarrow a=1$ and the restriction that only Voice plans cross; the composed
terminal target $s\mapsto b^*_{j(s)}(s)$ is strictly increasing on the flagged signal region, almost
surely on the flagged set; the inner pricing fixed point is unique; the flagged weight satisfies
$\Omega>0$; and the bidder's entry rule has the two properties
\begin{equation}
\label{eq:entry-props}
\Prb(\mathsf B=1\mid\mathcal I_H)=p(\mathcal I_H)
\qquad\text{and}\qquad
\mathsf B \perp v \mid \mathcal I_H .
\end{equation}
Then, on the flagged cell $\mathcal C_F$, the flagged tuple $\Sfl=(B^F,Q^F,a{=}1)$ makes the
pre-filing pooled history conditionally independent of $(v,s,\xi)$, and the flagged posterior, the
flagged price, the bidder's entry probability and the flagged cell average $M_F$ do not depend on
$\kappa$. In addition $\Omega$ does not depend on $\kappa$, so $\partial_\kappa\Omega=0$ at fixed
policies.
\end{lemma}

\begin{proof}
Write $\Xi:=(v,s,\xi)$ for the triple whose conditional independence is at issue, $z_{0:H}$ for the
noise vector, $\mathcal H_{f^-}^P$ for the pre-filing pooled history on a flagged path, and
$\mathcal I_H$ for the public information set at the control node, which the model sets equal to
$\Sfl$ on the flagged cell. The functions $u_1,u_2,u_3,u_5$ below are bounded and measurable and are
local to this proof.

\emph{Step 1 (where $\kappa$ lives).} At fixed cutoff and execution policies, and with no revision of
the path within the window, every policy object is a function of $(s;\tau,T)$ alone: the selected
plan $j=j(s)$, the stake path $B_j(s,\cdot)$, the order marks $q_{jd}(s)$, the terminal target
$b_j^*(s)$, the crossing date $c_j(s)$, the filing date $f_j(s)$, the disclosure indicator $D_j(s)$,
the stake at filing $B_j^F(s)$ and the flagged order $Q_j^F(s)$. Absence of feedback removes any
dependence on realised order flow or prices, and fixed policies remove any dependence on $\kappa$
through re-optimised cutoffs. Among the primitives, $\kappa$ appears in exactly one place, the law of
the ternary noise mark; the laws of $v$, $\varepsilon$ and $\xi$ and the remaining constants carry no
$\kappa$. Write $Q_\kappa$ for the law of $z_{0:H}$. At fixed policies, then, $\kappa$ enters the
model only through $Q_\kappa$.

\emph{Step 2 (the flagged set and its mass are free of $\kappa$).} The disclosure indicator
$D=\ind\{a=1,\ c(\tau)+T\le H\}$ is measurable, and by Step~1 it is a function of the signal alone.
Hence the flagged cell $\mathcal C_F=\{D=1\}$ lies in $\sigma(s)$ and is one fixed Borel subset of
the signal line, the same subset at every $\kappa$. The signal $s=v+\varepsilon$ is
$N(\mu_v,\sigma_v^2+\sigma_\varepsilon^2)$, a law free of $\kappa$ because the noise mark enters
neither $v$ nor $\varepsilon$. Therefore $\Omega=\Prb(s\in\mathcal C_F)$ is the same number at every
$\kappa$, which is the last conclusion of the statement.

\emph{Step 3 (the flagged tuple pins the signal, with a measurable inverse).} On $\mathcal C_F$ the
engagement flag is $a=1$, since disclosure implies engagement and only Voice plans cross, so the
third coordinate of $\Sfl$ is uninformative there. By the definition $Q_j^F=b_j^*(s)-B_j^F$,
\begin{equation}
\label{eq:flagged-target}
B^F_{j(s)}(s) \;+\; Q^F_{j(s)}(s) \;=\; b^*_{j(s)}(s).
\end{equation}
The composed terminal target on the right of \eqref{eq:flagged-target} is strictly increasing in the
signal on the flagged region and is therefore injective there. That monotonicity is consumed almost
surely on the flagged set, which is the only coherent reading when $B^F$ is continuum-valued: no
individual flagged tuple then carries positive probability, so injectivity holds almost surely on a
set of positive probability rather than on a set of atoms. Consequently the first two coordinates of
$\Sfl$ determine $b^*_{j(s)}(s)$, which determines $s$, which at fixed policies determines $j=j(s)$.
Write $\iota_F$ for the resulting inverse. It is strictly increasing on the image of the composed
target, and a monotone real function on a Borel set is Borel, so $\iota_F$ is measurable; no
measurable-selection theorem is needed. Since $\Omega>0$, the set on which this holds carries
strictly positive probability and the statement is not vacuous. Restricted to $\mathcal C_F$,
therefore, $\sigma(\Sfl)=\sigma(s)$ up to null sets. Strict monotonicity is what the step consumes:
if two signals produced the same flagged tuple while carrying different pooled paths,
\eqref{eq:flagged-target} would fail to be injective and the tuple would not pin the pooled path.

\emph{Step 4 (the pre-filing pooled history is a noise transform indexed by plan and signal).}
Define, for a plan $j$ and a signal $s$,
\begin{equation}
\label{eq:flagged-upsilon}
\Upsilon_{j,s}(z_{0:H}) \;:=\;
  \Bigl( \bigl(q_{j0}(s)+z_0,\dots,q_{j,f_j(s)-1}(s)+z_{f_j(s)-1}\bigr);\ \text{flag not landed by }
  f_j(s)-1 \Bigr),
\end{equation}
so that $\mathcal H_{f^-}^P=\Upsilon_{j(s),s}(z_{0:H})$ pointwise on $\mathcal C_F$. The indexing by
$(j,s)$ alone is exactly the absence of feedback: with re-optimisation inside the window the marks
would depend on realised flow, hence on past noise, and no such indexing would exist. The map
\eqref{eq:flagged-upsilon} is jointly measurable in $(s,z_{0:H})$. The filing date takes finitely
many values because the calendar is finite, and on each of its finitely many measurable level sets
the map is $(s,z)\mapsto(q_{j(s)d}(s)+z_d)_{d<t}$, measurable because
$q_{jd}=\Gamma\bigl(B_j(\cdot,d)-B_j(\cdot,d-1)\bigr)$ composes a Borel function, Voice stakes being
monotone in the signal and only Voice plans reaching $\mathcal C_F$, with the coarsening $\Gamma$,
itself Borel because it is ordered in the increment. This is unaffected by the size of the
blockholder's per-round order, which fixes the value the mark takes on a building round but not its
measurability in the signal. Patching finitely many measurable restrictions gives a measurable map.
The second component of \eqref{eq:flagged-upsilon} is a deterministic function of $s$ by Step~1 and
carries no information beyond $s$.

\emph{Step 5 (observed prices enlarge nothing).} Each pooled price is the unique fixed point of the
pricing map at its pooled history and is therefore a function of that history, and likewise the
flagged price is a function of the flagged information set. Adjoining the value of a function of a
generator of a $\sigma$-algebra to that $\sigma$-algebra does not enlarge it, so prices may be
ignored when the information content of $\mathcal I_H$ is computed.

\emph{Step 6 (an information sandwich).} On $\mathcal C_F$,
\begin{equation}
\label{eq:flagged-sandwich}
\sigma(s) \cap \mathcal C_F \;\subseteq\; \mathcal I_H \;\subseteq\;
  \sigma(s,z_{0:H}) \cap \mathcal C_F .
\end{equation}
For the left inclusion: the filing reveals $F=(B^F,a{=}1)$ and the flagged order is submitted and
priced, so $\Sfl$ belongs to $\mathcal I_H$; by Step~3 the tuple recovers $s$; and
$\mathcal C_F\in\sigma(s)$ by Step~2, so restricting to the flagged cell is itself an operation
within $\sigma(s)$. For the right inclusion: every component of $\mathcal I_H$, namely pooled order
flow, flag status, filing content, the flagged order and prices, is a measurable function of
$(s,z_{0:H})$ by Steps~1, 4 and~5. Everything below uses only \eqref{eq:flagged-sandwich}, so the
conclusion is robust to the exact content assigned to the control-node information set, provided it
satisfies the sandwich.

\emph{Step 7 (a freezing lemma).} Let $\mathcal G$ be a $\sigma$-algebra with $z_{0:H}$ independent
of $\mathcal G$, let $S$ be $\mathcal G$-measurable, let $w$ be bounded and jointly measurable, and
put $\bar w(x):=\int w(x,\zeta)\,Q_\kappa(\mathrm d\zeta)$. Then
$\Eop[w(S,z_{0:H})\mid\mathcal G]=\bar w(S)$. For a product form $w(x,\zeta)=u_1(x)u_2(\zeta)$ this
is $\Eop[u_1(S)u_2(z_{0:H})\mid\mathcal G]=u_1(S)\Eop[u_2(z_{0:H})]=\bar w(S)$, pulling out the
$\mathcal G$-measurable factor and using independence. The product forms generate the product
$\sigma$-algebra and are closed under multiplication; both sides of the asserted equality are linear
in $w$ and stable under bounded monotone limits, so the monotone class theorem extends the equality
to every bounded jointly measurable $w$.

\emph{Step 8 (the conditional independence).} By Step~3, conditioning on $\sigma(\Sfl)$ on the
flagged set is conditioning on $\sigma(s)$, so what has to be shown is
$\mathcal H_{f^-}^P \perp \Xi \mid s$ on $\mathcal C_F$. Given $s$, the middle coordinate of $\Xi$ is
degenerate, so it suffices to show, for bounded measurable $u_1$ on the pre-filing history space and
bounded measurable $u_2$ on $\mathbb R^2$,
\begin{equation}
\label{eq:flagged-factor}
\Eop\bigl[u_1(\mathcal H_{f^-}^P)\,u_2(v,\xi) \mid s\bigr]
 \;=\; \Eop\bigl[u_1(\mathcal H_{f^-}^P) \mid s\bigr]\cdot\Eop\bigl[u_2(v,\xi) \mid s\bigr]
 \qquad \text{a.s.\ on } \mathcal C_F,
\end{equation}
which is the standard characterisation of conditional independence given a $\sigma$-algebra. Put
$w(x,\zeta):=u_1\bigl(\Upsilon_{j(x),x}(\zeta)\bigr)$, bounded and jointly measurable by Step~4, so
that $u_1(\mathcal H_{f^-}^P)=w(s,z_{0:H})$. The noise vector is independent of
$(v,\varepsilon,\xi)$ and hence of $\Xi$, which is a measurable function of that triple. Apply
Step~7 with $\mathcal G=\sigma(v,s,\xi)$ and $S=s$:
$\Eop[w(s,z_{0:H})u_2(v,\xi)\mid v,s,\xi]=u_2(v,\xi)\,\bar w(s)$. Taking $\Eop[\,\cdot\mid s]$ of
both sides and pulling out the $\sigma(s)$-measurable factor $\bar w(s)$ gives
$\Eop[u_1u_2\mid s]=\bar w(s)\,\Eop[u_2\mid s]$; setting $u_2\equiv1$ in the same display gives
$\Eop[u_1\mid s]=\bar w(s)$, and substituting yields \eqref{eq:flagged-factor}. Unconditionally the
pre-filing pooled history depends strongly on the signal through the order marks, and the
conditioning event removes that dependence by pinning the signal, leaving the noise as the only
remaining randomness in \eqref{eq:flagged-upsilon}.

\emph{Step 9 (the flagged posterior).} Let $u_3$ be bounded, $\mathcal I_H$-measurable and supported
on $\mathcal C_F$. By the right inclusion of \eqref{eq:flagged-sandwich} and the Doob--Dynkin lemma,
$u_3=u_5(s,z_{0:H})$ for some bounded jointly measurable $u_5$. Running the computation of Step~8
with $u_5$ in place of $w$ gives $\Eop[u_2u_3]=\Eop\bigl[\bar u_5(s)\,\Eop[u_2\mid s]\bigr]$, while
the tower property together with $\Eop[u_3\mid s]=\bar u_5(s)$ from Step~7 gives
$\Eop\bigl[u_3\,\Eop[u_2\mid s]\bigr]=\Eop\bigl[\bar u_5(s)\,\Eop[u_2\mid s]\bigr]$. Hence
$\Eop[u_2u_3]=\Eop\bigl[\Eop[u_2\mid s]\,u_3\bigr]$ for every such $u_3$; since $\Eop[u_2\mid s]$ is
$\sigma(s)$-measurable and the left inclusion of \eqref{eq:flagged-sandwich} places
$\sigma(s)\cap\mathcal C_F$ inside $\mathcal I_H$, it is a version of $\Eop[u_2\mid\mathcal I_H]$ on
the flagged set:
\begin{equation}
\label{eq:flagged-posterior}
\Eop[u_2(v,\xi) \mid \mathcal I_H] \;=\; \Eop[u_2(v,\xi) \mid s]
 \qquad \text{a.s.\ on } \mathcal C_F .
\end{equation}
The extension from bounded to integrable $u_2$, needed at Step~11 for $u_2(v,\xi)=v$, follows by
applying \eqref{eq:flagged-posterior} to the truncations $\max(-n,\min(n,u_2))$ and passing to the
limit under conditional dominated convergence, the dominating variable being $|u_2|$, integrable
because $v$ is Gaussian. Separately, the engagement posterior satisfies $\pi(\mathcal I_H)=1$ on
$\mathcal C_F$, because the flagged cell lies inside $\{a=1\}$, and because the cell is itself
$\mathcal I_H$-measurable, being $\{D=1\}$ with $D$ part of the flag status carried in
$\mathcal I_H$.

\emph{Step 10 (the posterior does not move with $\kappa$).} By \eqref{eq:flagged-posterior} the
conditional law of $(v,\xi)$ given $\mathcal I_H$ on the flagged set equals its conditional law given
$s$. That law is $v\mid s \sim N\bigl(\mu_v+\beta(s-\mu_v),(1-\beta)\sigma_v^2\bigr)$ with
$\beta=\sigma_v^2/(\sigma_v^2+\sigma_\varepsilon^2)$, and $\xi\mid s \sim N(0,\sigma_\xi^2)$
independent of $v$. Neither depends on $\kappa$, which by Step~1 lives only in $Q_\kappa$ and enters
no coordinate of $\Xi$. With $\pi=1$ from Step~9, the flagged posterior over $(v,a,\xi)$ is a
function of the signal that does not move with $\kappa$, hence by Step~3 a function of $\Sfl$ that
does not move with $\kappa$.

\emph{Step 11 (the flagged price).} The flagged price solves the inner fixed point
$P=\Eop[Y\mid\mathcal I_H]$. On the flagged cell $a=1$ and $\pi=1$, so
$m_0+\Delta_m=m_1$. Only the two entry properties \eqref{eq:entry-props} are used. The entry rule of
the model has both, since the bidder's shock satisfies $\xi\mid\mathcal I_H \sim N(0,\sigma_\xi^2)$
independently of $v$ by Step~10 and the entry indicator is a function of $\xi$ and of
$\mathcal I_H$-measurable quantities; any other rule with these two properties works as well. Under
them,
\begin{equation}
\label{eq:flagged-innerFP}
P \;=\; \bigl(1-p(P)\bigr)\Bigl(\mu_v+\beta(s-\mu_v)+\Delta_V\Bigr) \;+\; p(P)\bigl(P+m_1\bigr),
\qquad
p(P) \;=\; 1-\Phi\!\left(\frac{P+K+m_1-\bar S}{\sigma_\xi}\right),
\end{equation}
where $\Eop[v\mid\mathcal I_H]=\Eop[v\mid s]=\mu_v+\beta(s-\mu_v)$ by
\eqref{eq:flagged-posterior}. The data of \eqref{eq:flagged-innerFP} are the single number
$\Eop[v\mid s]$ and the primitive parameters, none of which depends on $\kappa$ by Step~10.
Uniqueness of the inner fixed point makes the solution a function of the data alone, so equal data
give equal solutions. Uniqueness is load-bearing rather than decorative: without it an extraneous
selection among fixed points could vary with $\kappa$ even though the map does not. Hence the flagged
price is a function of the signal, and by Step~3 of $\Sfl$, that does not move with $\kappa$.

\emph{Step 12 (the entry probability).} With $\pi=1$ from Step~9, the flagged entry probability is
$p=1-\Phi\bigl((P^F+K+m_1-\bar S)/\sigma_\xi\bigr)$ on the flagged cell. Every argument is free of
$\kappa$, by Step~11 for the price and by the primitive table for the constants. Write $p_F(s)$ for
it. It is measurable as a composition of the monotone map $s\mapsto\Eop[v\mid s]$, the fixed-point
solution and the normal distribution function, and it takes values in $(0,1)$.

\emph{Step 13 (the flagged cell average).} The kernel of the premium is $h=\pi p$, so $h=p_F(s)$ on
the flagged cell by Steps~9 and~12. Since $\Omega>0$, the elementary definition of conditioning on
a positive-probability event gives
\begin{equation}
\label{eq:flagged-MF}
M_F \;=\; \Delta_m\,\Eop[h \mid D=1]
     \;=\; \frac{\Delta_m\,\Eop\bigl[p_F(s)\,\ind\{s \in \mathcal C_F\}\bigr]}{\Omega} .
\end{equation}
The numerator of \eqref{eq:flagged-MF} integrates a bounded measurable function free of $\kappa$
(Step~12) over a Borel set free of $\kappa$ (Step~2) against the law of the signal, itself free of
$\kappa$ (Step~2), and the denominator is free of $\kappa$ and strictly positive. Hence $M_F$ does
not move with $\kappa$, which completes the four invariance conclusions.
\end{proof}

Two boundaries of the result are part of it. Nothing above asserts that the pre-filing pooled price
does not move with $\kappa$: that price conditions on order flow whose law is $Q_\kappa$, and it
moves with $\kappa$ in general. With $J=P^F-P_{\ND}$, where $P_{\ND}$ is the last pre-filing pooled
price at the same realised order flow, and with $\partial_\kappa P^F=0$ on the flagged set, wherever
the derivative exists $\partial_\kappa J=-\partial_\kappa P_{f^-}^P$. The result therefore makes no
claim about the filing-day jump. And the result is a fixed-policy statement: if cutoffs or execution
paths move with $\kappa$, Step~1 is unavailable and the conclusions do not follow.

%==============================================================================
% The partition and the factorisation  S = (1 - Omega) S_P
%==============================================================================

\providecommand{\Prb}{\mathbb{P}}
\providecommand{\Eop}{\mathbb{E}}
\providecommand{\ind}{\mathbf{1}}
\providecommand{\Dact}{\Delta_{\mathrm{act}}}

\subsection{The partition and the factorisation}
\label{sec:partition-factorisation}

The disclosure rule $(\tau,T)$ acts on the model first as a partition of histories and only then as
a statement about prices. This subsection separates the two. The partition is combinatorial and
needs no restriction on cell mass; the two-cell decomposition of the premium needs an interior
partition; the factorisation of the liquidity sensitivity needs, in addition, that the flagged
endpoint and the flagged weight are both free of $\kappa$.

\subsubsection*{Standing conditions}

The following are in force throughout the subsection.

\begin{enumerate}[label=(S\arabic*),leftmargin=3.0em,itemsep=2pt]
\item \label{p:space} \emph{One probability space.} The primitive vector, the value $v$, the signal
  noise $\varepsilon$, the bidder draw $\xi$ and the noise-trader marks $z_{0:H}$, has a joint law
  on a finite product of Polish spaces, with the noise marks valued in the finite set
  $\{-\bar z,0,+\bar z\}$. The signal is $s = v + \varepsilon$.
\item \label{p:menu} \emph{A finite menu and a finite calendar.} The plan menu $\mathcal J$ is
  finite, the horizon $H$ is finite, and the window margin satisfies $T \in \{1,\dots,H\}$.
\item \label{p:cutoff} \emph{A cutoff selection map.} The signal is carried into a plan by a Borel step
  function $j(\cdot)$ with breakpoints $k_1 \le \dots \le k_{J-1}$.
\item \label{p:paths} \emph{Monotone Voice paths and a clean start.} Each plan $j$ carries a Borel stake
  path $B_j(s,\cdot)$ with $B_j(s,-1) = b_0 < \tau$; for a Voice plan the path is weakly increasing
  in the signal and weakly increasing in the calendar date.
\item \label{p:clock} \emph{Legal-clock discipline.} Only Voice plans cross the threshold; the
  crossing date $c_j(s)$ is the first date at which the path reaches $\tau$, or $+\infty$ if the
  path never reaches it; the filing lands exactly at $f_j = c_j + T$ and only through the disclosure
  node; the filing reports the stake truthfully.
\item \label{p:public} \emph{The flag is public.} Every control-node history carries the public
  coordinate ``the filing has landed by date $d$'', so the control-node information set
  $\mathcal I_H$ contains that coordinate at $d = H$ on both cells. On the pooled cell
  $\mathcal I_H$ also contains the pooled history up to the control node.
\item \label{p:nofeedback} \emph{No-feedback timing.} The executed path, the order marks, the
  terminal target, the crossing date, the filing date, the filing stake and the flagged order are
  Borel functions of the plan and the signal alone, with the filing stake and flagged order defined
  on the flagged cell. Neither realised order flow nor a realised price enters any of them.
\item \label{p:kernel} \emph{A bounded, pinned kernel.} The engagement-premium kernel is
  $h(\mathcal I) = \pi(\mathcal I)\,p(\mathcal I)$, where the engagement posterior
  $\pi(\mathcal I) = \Prb(a = 1 \mid \mathcal I)$ takes values in $[0,1]$ and the bidder-entry
  probability $p(\mathcal I)$ takes values in $(0,1)$ and is a continuous function of the posterior
  and the price. The pricing rule pins one version of $\Eop[Y \mid \mathcal I]$ rather than an
  equivalence class.
\item \label{p:wedge} \emph{A finite wedge.} The premium wedge $\Delta_m$ is a finite, strictly
  positive constant, and $\Dact = \Delta_m\,\Eop[h(\mathcal I_H)]$.
\item \label{p:kappa} \emph{Liquidity enters in one place.} The intensity $\kappa$ appears in the
  primitives only in the law of the ternary noise mark. The laws of $v$, $\varepsilon$ and $\xi$
  and the remaining constants carry no $\kappa$.
\item \label{p:fixed} \emph{Fixed policies.} The plan menu, the execution policies and the cutoff
  vector $k$ are held fixed in $\kappa$ and held fixed across any two rules compared.
For the remaining conditions, let $\theta=(j(s),s)$ denote the blockholder's type,
$m_d(\theta)=q_{j(s)d}(s)$ its order mark at date $d$, and $r$ an arbitrary disclosure rule under
comparison. Write $\mathsf S_{F,r}=(B_r^F,Q_r^F,a{=}1)$ for that rule's flagged tuple.
\item \label{p:noise} \emph{A common noise channel.} The vector $z_{0:H}$ is independent of
  $(v,\varepsilon,\xi)$ and has the common rule-invariant law $Q_\kappa$. Pooled flow is
  $X_d=m_d(\theta)+z_d$. In the order-size-two model,
  $m_d(\theta)\in\{0,2\bar z\}$, and the noise marks are independent across dates and have
  probabilities $(\kappa/2,1-\kappa,\kappa/2)$ on
  $\{-\bar z,0,+\bar z\}$.
\item \label{p:control} \emph{The control-node convention.} On the flagged cell the information set
  is exactly $\sigma(\mathsf S_{F,r})$, up to pinned price coordinates. On the pooled cell it is
  exactly the displayed public pooled history, again up to pinned prices. At each depth $d$, the
  same convention applies to the cells $\mathcal C_{F,r}^d$ and $\mathcal C_{P,r}^d$ defined by
  whether the filing has landed by $d$.
\item \label{p:decoder} \emph{Identified flagged types.} On rule $r$'s flagged signal region, the
  composed terminal target $g(s)=b^*_{j(s)}(s)$ is strictly increasing and
  $B^F_r+Q^F_r=g(s)$. Pointwise strict increase is imposed for pointwise experiment comparisons;
  for posterior conclusions, strict increase outside a prior-null set is enough.
\end{enumerate}

\subsubsection*{The partition}

\begin{lemma}[Disclosure partition]
\label{lem:partition}
Under \ref{p:space}--\ref{p:public} and \ref{p:control}, the disclosure indicator
$D = \ind\{a = 1,\ c(\tau) + T \le H\}$ is a measurable function of the control-node history, and
the flagged cell $\mathcal C_F = \{D = 1\}$ and the pooled cell $\mathcal C_P = \{D = 0\}$ are
disjoint and cover the space. Consequently the flagged weight $\Omega = \Prb(D = 1)$ is defined and
$\Prb(D = 0) = 1 - \Omega$. No restriction on the mass of either cell is used.
\end{lemma}

\begin{proof}
\emph{Step 1 (the selection map is Borel).} By \ref{p:space} the signal $s = v + \varepsilon$ is a
sum of two coordinates of a Borel vector and is therefore Borel. Condition \ref{p:cutoff} makes
$j(\cdot)$ Borel. Its step form has finitely many breakpoints by \ref{p:menu}. The set
$\{s : j(s) = j\}$ is Borel for each
$j \in \mathcal J$.

\emph{Step 2 (the engagement flag is measurable).} The realised flag is $a = a_{j(s)}$, the
composition of the Borel map of Step~1 with the map $j \mapsto a_j$ defined on the finite set
$\mathcal J$. A map on a finite set carrying its discrete $\sigma$-algebra is measurable, so
$\{a = 1\}$ is Borel and lies in $\sigma(s)$.

\emph{Step 3 (the crossing date is attained, unique and measurable).} Fix a Voice plan $j$. By
\ref{p:paths} the map $s \mapsto B_j(s,d)$ is weakly increasing, and the preimage of
$[\tau,\infty)$ under a weakly monotone real function is an interval, so
$\{s : B_j(s,d) \ge \tau\}$ is Borel at each date $d$. The calendar $\{0,\dots,H\}$ is finite by
\ref{p:menu}, so
\[
  \{c_j \le d\} \;=\; \bigcup_{d' \le d}\{s : B_j(s,d') \ge \tau\}
\]
is a finite union of Borel sets and $c_j$ is a Borel map into $\{0,\dots,H\} \cup \{+\infty\}$. Its
value is unique: a nonempty subset of the finite well-ordered set $\{0,\dots,H\}$ has exactly one
minimum, and by \ref{p:clock} the crossing date is that minimum; if the subset is empty then
$c_j = +\infty$.

\emph{Step 4 (the disclosure indicator is measurable).} By \ref{p:clock} a filing requires
$a = 1$, and $f_j = c_j + T \le H$ holds exactly when $c_j \le H - T$, the convention
$(+\infty) + T = +\infty > H$ being available because $H$ is finite by \ref{p:menu}. Hence
\begin{equation}
\label{eq:cell-union}
  \{D = 1\} \;=\; \bigcup_{j \in \mathcal J\,:\,a_j = 1}
    \Bigl( \{s : j(s) = j\} \;\cap\; \{s : c_j(s) \le H - T\} \Bigr).
\end{equation}
Each set on the right is Borel by Steps~1 and~3, and the union is finite by \ref{p:menu}, so
$\{D = 1\}$ is Borel and $D$ is a measurable $\{0,1\}$-valued map. Record one fact for later: at a
fixed cutoff policy every ingredient of \eqref{eq:cell-union} is a function of $s$ alone, so
$\{D = 1\} \in \sigma(s)$. Record a second: regularity of $s \mapsto B_j(s,d)$ was needed only for
Voice plans, the conjunct $a = 1$ having removed the others from the union.

\emph{Step 5 (exactly one cell, at the level of states).} By Step~4, $D$ is the indicator of a Borel
event, so at every sample point it takes exactly one of the values $0$ and $1$. Therefore
$\mathcal C_F \cap \mathcal C_P = \varnothing$ and $\mathcal C_F \cup \mathcal C_P$ is the whole
space.

\emph{Step 6 (the cell is a function of the observed history).} Step~5 partitions the space of
states; the claim is about control-node histories. By \ref{p:public} every control-node history carries the
public coordinate ``the filing has landed by date $d$'', and
$\mathcal I_H$ contains that coordinate at $d = H$ on both cells. Under \ref{p:control}, flagged
tuples and pooled histories occupy disjoint tagged output spaces. By \ref{p:clock} the filing lands exactly at
$c + T$ and only through the disclosure node, so the coordinate ``the filing has landed by $H$''
equals $\ind\{f \le H\}$; combining this with $D = \ind\{a = 1,\ c < \infty,\ f \le H\}$ and with the
restriction that only Voice plans cross, that coordinate equals $D$. Hence $D$ is measurable with
respect to the control-node history and the map from a history to its cell is single-valued.
When both cells have positive probability, the map is onto
$\{\mathcal C_F, \mathcal C_P\}$. When $\Omega \in \{0,1\}$, one cell is null and the map
hits only the other almost surely, as Step~7 records. The public flag coordinate is what carries
this step: without it
$D$ would still be a measurable function of the state by Step~4, while the cell assignment would
stop being a function of what is observed at the control node.

\emph{Step 7 (the weight).} By Steps~5 and~6 the two cells are measurable, disjoint and exhaustive,
so $\Omega = \Prb(D = 1)$ is defined and $\Prb(D = 0) = 1 - \Omega$. Steps~1 to~7 used no
restriction on the value of $\Omega$, which may be $0$ or $1$.
\end{proof}

\subsubsection*{The two-cell decomposition}

\begin{lemma}[Two-cell decomposition of the premium]
\label{lem:two-cell}
Under \ref{p:space}--\ref{p:wedge}, with the partition of Lemma~\ref{lem:partition}, and whenever
$0 < \Omega < 1$,
\begin{equation}
\label{eq:two-cell}
  \Dact \;=\; \Omega\,M_F \;+\; (1 - \Omega)\,M_P,
  \qquad
  M_F = \Delta_m\Eop[h \mid D = 1],
  \quad
  M_P = \Delta_m\Eop[h \mid D = 0].
\end{equation}
At $\Omega = 1$ the identity reads $\Dact = M_F$ and at $\Omega = 0$ it reads $\Dact = M_P$; in each
of those two cases the average over the null cell is undefined rather than zero, and it is not
determined by the law. That $M_P$ is defined as a number, rather than an equivalence class, rests
on \ref{p:kernel} pinning a version of the price; Borel regularity of the stake paths in
\ref{p:paths} does not by itself select a version of a conditional expectation.
\end{lemma}

\begin{proof}
\emph{Step 1 (the kernel is bounded).} By \ref{p:kernel}, $\pi \in [0,1]$ and $p \in (0,1)$, so
their product satisfies $0 \le h \le 1$ pointwise.

\emph{Step 2 (the kernel is a random variable).} The posterior
$\pi(\mathcal I_H) = \Prb(a = 1 \mid \mathcal I_H)$ is a version of a conditional probability and is
$\mathcal I_H$-measurable, and so is the price $\Eop[Y \mid \mathcal I_H]$, whose version is pinned
by \ref{p:kernel}. That is what makes $h(\mathcal I_H)$ one well-defined random variable rather than
an equivalence class from which the decomposition would have to choose. The entry probability is the
composition of the posterior and the price with a continuous map, hence $\mathcal I_H$-measurable,
and a product of measurable functions is measurable. With Step~1 the kernel is integrable.

\emph{Step 3 (the pointwise split).} By Lemma~\ref{lem:partition} exactly one of $\ind\{D = 1\}$ and
$\ind\{D = 0\}$ equals one at every sample point and the other equals zero, so
$h = h\,\ind\{D = 1\} + h\,\ind\{D = 0\}$ pointwise.

\emph{Step 4 (integration).} Each summand in Step~3 is bounded in absolute value by $1$ by Step~1,
hence integrable on the single probability space of \ref{p:space}. Linearity of the integral gives
$\Eop[h] = \Eop[h\,\ind\{D = 1\}] + \Eop[h\,\ind\{D = 0\}]$.

\emph{Step 5 (the interior case).} Suppose $0 < \Omega < 1$. Both conditioning events have strictly
positive probability, so the elementary conditional expectations
$\Eop[h \mid D = 1] := \Eop[h\,\ind\{D = 1\}]/\Omega$ and
$\Eop[h \mid D = 0] := \Eop[h\,\ind\{D = 0\}]/(1 - \Omega)$ are defined, and Step~1 bounds both by
$1$. This is the defining property of conditioning on an event of positive probability. Substituting
into Step~4 gives $\Eop[h] = \Omega\,\Eop[h \mid D = 1] + (1 - \Omega)\,\Eop[h \mid D = 0]$.

\emph{Step 6 (scaling to premia).} By \ref{p:wedge} the wedge $\Delta_m$ is finite and strictly
positive. Multiplying Step~5 by $\Delta_m$ and reading off $\Dact$, $M_F$ and $M_P$ gives
\eqref{eq:two-cell}.

\emph{Step 7 (the case $\Omega = 1$).} Then $\Prb(D = 0) = 0$, and by Step~1
$|\Eop[h\,\ind\{D = 0\}]| \le \Prb(D = 0) = 0$, so that term of Step~4 vanishes and
$\Eop[h] = \Eop[h\,\ind\{D = 1\}] = \Omega\,\Eop[h \mid D = 1] = \Eop[h \mid D = 1]$, the middle
equality being Step~5's definition, available because $\Omega = 1 > 0$. Multiplying by $\Delta_m$
gives $\Dact = M_F$.

\emph{Step 8 (at $\Omega = 1$ the pooled average is not determined by the law).} Step~5's definition
does not apply to $M_P$, since it would divide by $\Prb(D = 0) = 0$. Conditioning on the generated
$\sigma$-algebra supplies no value either. Let $\mathcal D := \sigma(D)$, whose atoms are
$\{D = 1\}$ and $\{D = 0\}$, and let $x$ be any real number. The function
\begin{equation}
\label{eq:version}
  \Delta_m\,\Eop[h \mid D = 1]\,\ind\{D = 1\} \;+\; x\,\ind\{D = 0\}
\end{equation}
is $\mathcal D$-measurable and integrates to $\Delta_m\Eop[h\,\ind_A]$ over every
$A \in \mathcal D$, because two such functions with different values of $x$ differ only on the null
set $\{D = 0\}$. Every real $x$ therefore yields a version of $\Delta_m\Eop[h \mid \mathcal D]$, so
no value of $M_P$ is determined by the law, and the statement at $\Omega = 1$ is Step~7's one-term
equation.

\emph{Step 9 (the case $\Omega = 0$).} Mirror Steps~7 and~8 with the cells exchanged:
$|\Eop[h\,\ind\{D = 1\}]| \le \Prb(D = 1) = 0$, so $\Eop[h] = \Eop[h \mid D = 0]$ and
$\Dact = M_P$, while $M_F$ is left undetermined by the argument of Step~8.
\end{proof}

\medskip
\noindent\emph{The identity is tied to this partition.} The identity \eqref{eq:two-cell} is tied to the disclosure partition and to no other. The flagged
weight factors as $\Omega = \Prb(a = 1)\,\omega_a$ with $\omega_a = \Prb(D = 1 \mid a = 1)$, so
$\Omega \le \Prb(a = 1)$, strictly when $\Prb(a = 1) > 0$ and $\omega_a < 1$. Averaging the kernel
over $\{a = 1\}$ and $\{a = 0\}$ therefore gives a different decomposition, with the engaged but
unflagged histories sitting in the pooled cell of \eqref{eq:two-cell}.
\medskip

\subsubsection*{The factorisation}

Write $\mathcal S = |\partial_\kappa \Dact|$ for the liquidity sensitivity of the premium and
$\mathcal S_P = |\partial_\kappa M_P|$ for the liquidity sensitivity of the pooled cell.

\begin{proposition}[Factorisation of the liquidity sensitivity]
\label{prop:factorisation}
Assume \ref{p:space}--\ref{p:fixed}, the partition of Lemma~\ref{lem:partition}, the two-cell
decomposition of Lemma~\ref{lem:two-cell}, an interior partition $0 < \Omega < 1$ at the policy under
consideration, the invariance of the flagged endpoint in $\kappa$ at fixed policies, and, as an
explicit hypothesis, differentiability of $\kappa \mapsto M_P$. Then
\begin{equation}
\label{eq:factorisation}
  \mathcal S(\kappa,\tau,T) \;=\; \bigl(1 - \Omega(\tau,T)\bigr)\,\mathcal S_P(\kappa,\tau,T),
\end{equation}
exactly. The same factor scales the total variation of $\Dact$ over any grid in $\kappa$, and there
no differentiability is used.
\end{proposition}

\begin{proof}
\emph{Step 1 (the two-cell identity).} The partition is interior at the policy considered, so
Lemma~\ref{lem:two-cell} applies in its non-degenerate branch and \eqref{eq:two-cell} holds at
$(\kappa,\tau,T)$, with $M_F$ and $M_P$ the two cell averages.

% The flagged-cell invariance is the companion result of this subsection ("the flagged cell is
% kappa-free"); the cross-reference is wired when the appendix is assembled.
\emph{Step 2 (the flagged cell does not move with $\kappa$).} At the fixed cutoff and execution
policies of \ref{p:fixed}, the flagged posterior, the flagged price, the entry probability and $M_F$
are invariant to $\kappa$, by the flagged-cell invariance assumed in the statement. Hence
$\partial_\kappa M_F = 0$, and indeed $M_F$ is one and the same number at any two values of
$\kappa$.

\emph{Step 3 (the flagged weight does not move with $\kappa$ either).} This is derived, not assumed.
By Step~4 of Lemma~\ref{lem:partition}, $\{D = 1\}$ is the finite union \eqref{eq:cell-union}, each
of whose ingredients is a function of the signal alone once the cutoff policy is fixed. By
\ref{p:nofeedback} the executed path $B_j(s,d)$ is a function of $(j,s,d)$ alone, so no realised
order flow and no realised price enters it, and with the policies frozen by \ref{p:fixed} the
indicator $D$ is a function of $s$ alone. The law of $s$ carries no $\kappa$: by \ref{p:space} the
signal is $v + \varepsilon$, and by \ref{p:kappa} the laws of $v$ and $\varepsilon$ do not involve
$\kappa$. Indeed $\kappa$ enters the primitives in one place only, the law of the noise mark $z_d$,
which reaches observed pooled order flow $X_d = q_{jd} + z_d$ and reaches nothing that $D$ depends
on. Hence $\Omega = \Prb(D = 1)$ is literally the same number at every $\kappa$, and in particular
$\partial_\kappa\Omega = 0$. This step consumes the freezing of the cutoff vector in \ref{p:fixed},
and no property of the disclosure rule beyond the partition.

\emph{Step 4 (differentiating in $\kappa$).} Two of the three factors on the right of
\eqref{eq:two-cell} are constant in $\kappa$: $M_F$ by Step~2 and $\Omega$ by Step~3. Neither of
those steps uses the present one, so nothing here is circular. With those factors constant,
$\kappa \mapsto \Dact$ is an affine image of $\kappa \mapsto M_P$, which the statement's
differentiability hypothesis makes differentiable; hence $\partial_\kappa\Dact$ exists and equals
$(1 - \Omega)\,\partial_\kappa M_P$. Written with one term per factor that could carry $\kappa$, the
same computation reads
\begin{equation}
\label{eq:three-term}
  \partial_\kappa\Dact
  \;=\; \Omega\,\partial_\kappa M_F
  \;+\; (1 - \Omega)\,\partial_\kappa M_P
  \;+\; (\partial_\kappa\Omega)\,(M_F - M_P),
\end{equation}
the three terms being the flagged cell's own motion, the pooled cell's own motion, and the
reallocation of mass between the cells. Step~2 kills the first and Step~3 the third. The third
deletion assumes nothing about $M_F - M_P$, which is neither small nor signed here; the term goes
because its coefficient is zero.

\emph{Step 5 (the factorisation).} Since $0 < \Omega < 1$ we have $1 - \Omega \in (0,1)$ and
$|1 - \Omega| = 1 - \Omega$. Taking absolute values in the conclusion of Step~4 gives
\eqref{eq:factorisation}, exactly. Step~3 licenses writing the argument of $\Omega$ as $(\tau,T)$
rather than $(\kappa,\tau,T)$, and that convention is kept from here on. Two consequences are worth
recording. First, $\mathcal S > 0$ exactly when $\mathcal S_P > 0$. Second, when $\mathcal S_P > 0$
the map $x \mapsto |x|$ is differentiable at $\partial_\kappa M_P \ne 0$, so $\mathcal S_P$ inherits
whatever differentiability $\partial_\kappa M_P$ has. That transfer creates
no differentiability; it moves it.

\emph{Step 6 (the aggregate form, with no differentiability used).} Fix any grid
$\kappa_0 < \kappa_1 < \dots < \kappa_n$ inside the maintained parameter set, with the partition
interior at each node. Applying Step~1 at $\kappa_i$ and at $\kappa_{i+1}$, and using Steps~2 and~3
in their constancy form, that $M_F$ and $\Omega$ are one and the same number at the two nodes, which
a vanishing derivative at a single $\kappa$ would not deliver, gives
$\Dact(\kappa_{i+1}) - \Dact(\kappa_i) = (1 - \Omega)\bigl[M_P(\kappa_{i+1}) - M_P(\kappa_i)\bigr]$.
Summing absolute values,
\begin{equation}
\label{eq:tv-form}
  \sum_{i=0}^{n-1}\bigl\lvert \Dact(\kappa_{i+1}) - \Dact(\kappa_i)\bigr\rvert
  \;=\; (1 - \Omega)\sum_{i=0}^{n-1}\bigl\lvert M_P(\kappa_{i+1}) - M_P(\kappa_i)\bigr\rvert,
\end{equation}
which is the total-variation clause, finite differences throughout. The same argument runs for any
aggregator of the increment vector that is positively homogeneous of degree one, such as total
variation, mean absolute slope, or the supremum of the absolute increments, because $1 - \Omega$ is
one nonnegative scalar common to every increment. A measurement convention that reports
$\kappa$-sensitivity as a total variation over a grid rather than as a pointwise derivative is then
measuring the same functional.
\end{proof}

% ===========================================================================
% The weight leg at the threshold margin: a tighter threshold raises the
% flagged weight, so the pooled share weakly falls.
% ===========================================================================

\providecommand{\Prb}{\mathbb{P}}
\providecommand{\ind}{\mathbf{1}}

\begin{lemma}[The weight leg at the threshold margin]
\label{lem:threshold-weight}
Fix the liquidity $\kappa$, the window margin $T$, and the plan and cutoff policies. Compare two
threshold margins with $b_0 < \tau' < \tau$, so that $\tau'$ is the tighter threshold, and assume
$\Omega(\tau',T) < 1$. The configuration $b_0 < \tau'$ is what puts the clock equivalence in force
at both thresholds, together with Voice date-monotonicity (standing condition \ref{p:paths}) and
the restriction that only Voice plans cross (standing condition \ref{p:clock}). Then
\begin{enumerate}[label=(\roman*),leftmargin=2.4em,itemsep=3pt]
\item the flagged cells are nested, $\mathcal C_F(\tau,T) \subseteq \mathcal C_F(\tau',T)$, and
every history the tighter threshold newly flags carries the engagement flag $a=1$;
\item the flagged weight rises, $\Omega(\tau',T) \ge \Omega(\tau,T)$, so the weight leg
\begin{equation}
\label{eq:threshold-weight-leg}
W_\tau \;=\; \frac{1-\Omega(\tau',T)}{1-\Omega(\tau,T)}
\end{equation}
is defined and attenuates: $0 < W_\tau \le 1$;
\item if in addition $\Omega(\tau,T) > 0$, $\mathcal S_P(\kappa,\tau,T) > 0$, and the hypotheses of
Proposition~\ref{prop:factorisation} hold at both thresholds, the sensitivity
ratio at this margin factors into the weight leg and the composition leg
$C_\tau = \mathcal S_P(\kappa,\tau',T)/\mathcal S_P(\kappa,\tau,T)$,
\begin{equation}
\label{eq:threshold-ratio}
\frac{\mathcal S(\kappa,\tau',T)}{\mathcal S(\kappa,\tau,T)} \;=\; W_\tau\,C_\tau .
\end{equation}
\end{enumerate}
\end{lemma}

\begin{proof}
Fixed policies hold the map from the signal to the plan, the executed stake path of each plan, and
the law of the signal common to the two rules. The disclosure rule writes the stake path with no
threshold argument, while the crossing date, the filing date, the stake at filing and the
disclosure indicator each carry one. The only objects that move when the threshold moves are
therefore those four, and every comparison below is made at one and the same plan, one and the same
path, and one and the same signal law.

\emph{The product form of the disclosure indicator.} For every plan $j$, every signal $s$, and each
of the two thresholds,
\begin{equation}
\label{eq:threshold-product}
D_j(s;\tau,T) \;=\; \ind\{a_j = 1\}\cdot\ind\{B_j(s,H-T) \ge \tau\} .
\end{equation}
If $a_j = 0$ both sides vanish, the left side because engagement is a conjunct of disclosure and the
right side through its first factor. If $a_j = 1$ the left side is the event that the filing lands
by the horizon: the filing date is the crossing date plus the window margin, and with a finite
window margin the filing landing by the horizon already forces the crossing to be finite, so the
requirement that a crossing occur adds nothing. The clock equivalence then converts the filing
landing by the horizon into the stake $T$ rounds before the horizon having reached the threshold,
which is the second factor. Because $b_0 < \tau' < \tau$, the crossing is the first passage at each
of the two thresholds and the equivalence is available at both.

\emph{The inclusion.} Take any control-node history that the looser threshold flags,
$D_j(s;\tau,T) = 1$. By \eqref{eq:threshold-product} this says $a_j = 1$ and
$B_j(s,H-T) \ge \tau$. Since $\tau' < \tau$, the same real number satisfies
$B_j(s,H-T) \ge \tau > \tau'$, and reading \eqref{eq:threshold-product} at $\tau'$ for the same plan
and the same signal gives $D_j(s;\tau',T) = 1$. The number $B_j(s,H-T)$ that appears in the two
comparisons is one number, and this is where fixed policies are consumed: were the path free to move
with the threshold, the comparison at $\tau'$ would be about a different path. The history was
arbitrary, so $\mathcal C_F(\tau,T) \subseteq \mathcal C_F(\tau',T)$, and taking complements in the
exclusive and exhaustive partition gives
$\mathcal C_P(\tau',T) \subseteq \mathcal C_P(\tau,T)$: a tighter threshold can only shrink the
pooled cell. Every newly flagged history satisfies $D(\tau') = 1$, and by
\eqref{eq:threshold-product} that carries $a = 1$ as a conjunct, so the engagement flag is
identically one on the newly flagged set. Nestedness is a conclusion here rather than a hypothesis.
This is clause~(i).

\emph{The weight leg.} The joint law of the primitives is itself a primitive and does not depend on
the threshold, which enters only through the flagged event. Monotonicity of a probability measure
applied to the inclusion gives
\[
\Omega(\tau',T) \;=\; \Prb\bigl(\mathcal C_F(\tau',T)\bigr)
\;\ge\; \Prb\bigl(\mathcal C_F(\tau,T)\bigr) \;=\; \Omega(\tau,T),
\]
so $1-\Omega(\tau',T) \le 1-\Omega(\tau,T)$. The hypothesis $\Omega(\tau',T) < 1$ makes the
numerator of \eqref{eq:threshold-weight-leg} strictly positive, and the same hypothesis with the
inequality just established gives $\Omega(\tau,T) \le \Omega(\tau',T) < 1$, so the denominator is
strictly positive as well. The quotient is therefore defined, and dividing a weakly smaller strictly
positive number by a strictly positive one gives $0 < W_\tau \le 1$. This is clause~(ii).

\emph{The ratio identity.} With $\Omega(\tau,T) > 0$ the weight-leg inequality puts
both flagged weights strictly inside the unit interval, so Proposition~\ref{prop:factorisation}
applies at each threshold and the noise sensitivity of the expected engagement premium is the
pooled share times the pooled sensitivity,
\[
\mathcal S(\kappa,\tau,T) \;=\; \bigl(1-\Omega(\tau,T)\bigr)\,\mathcal S_P(\kappa,\tau,T),
\qquad
\mathcal S(\kappa,\tau',T) \;=\; \bigl(1-\Omega(\tau',T)\bigr)\,\mathcal S_P(\kappa,\tau',T) .
\]
The flagged weight does not move with $\kappa$, which is what licenses writing it as
$\Omega(\cdot,T)$ with no liquidity argument on either side. Since $1-\Omega(\tau,T) > 0$ and
$\mathcal S_P(\kappa,\tau,T) > 0$, the denominator $\mathcal S(\kappa,\tau,T)$ is strictly positive
and the quotient of the two displays is defined. Dividing and grouping the two factors gives
\eqref{eq:threshold-ratio}, which is clause~(iii).
\end{proof}

\noindent
The weight leg at this margin is signed outright, and it is signed by the geometry of the flagged
set alone: no restriction on the pooled experiment and no restriction on the noise-trading intensity
enters clauses~(i) and~(ii). What clause~(iii) then supplies is the exact form in which the
composition leg has to be signed: $W_\tau C_\tau \le 1$ if and only if $C_\tau \le 1/W_\tau$, and
by clause~(ii) that ceiling is at least one.

% Clock theorem: a shorter clock lowers noise sensitivity iff W_T C_T is at most one.

\begin{theorem}[The clock dial]
\label{thm:clock}
Fix the liquidity $\kappa$, the threshold margin $\tau$, and the plan and cutoff policies. Compare
two window margins $T' < T$, so that $T'$ is the shorter clock. Assume $b_0 < \tau$, that Voice
stake paths are weakly increasing in the calendar date (standing condition \ref{p:paths}), and
that only Voice plans cross (standing condition \ref{p:clock}), which is the configuration that
puts the clock equivalence in force at both windows. Assume $0 < \Omega(\tau,T) < 1$ and
$0 < \Omega(\tau,T') < 1$, and $\mathcal S_P(\kappa,\tau,T) > 0$. Assume that the hypotheses of
Proposition~\ref{prop:factorisation} hold at both clocks. Write
\begin{equation}
\label{eq:clock-legs}
W_T \;=\; \frac{1-\Omega(\tau,T')}{1-\Omega(\tau,T)},
\qquad
C_T \;=\; \frac{\mathcal S_P(\kappa,\tau,T')}{\mathcal S_P(\kappa,\tau,T)} .
\end{equation}
Then
\begin{enumerate}[label=(\roman*),leftmargin=2.4em,itemsep=3pt]
\item the weight leg attenuates: $0 < W_T \le 1$;
\item the sensitivity ratio factors into the two legs,
\begin{equation}
\label{eq:clock-ratio}
\frac{\mathcal S(\kappa,\tau,T')}{\mathcal S(\kappa,\tau,T)} \;=\; W_T\,C_T ;
\end{equation}
\item the shorter clock lowers noise sensitivity exactly when the product of the two legs is at
most one,
\begin{equation}
\label{eq:clock-iff}
\mathcal S(\kappa,\tau,T') \;\le\; \mathcal S(\kappa,\tau,T)
\qquad\Longleftrightarrow\qquad
W_T\,C_T \;\le\; 1 .
\end{equation}
\end{enumerate}
\end{theorem}

\begin{proof}
Fixed policies hold the map from the signal to the plan, the executed stake path of each plan, and
the law of the signal common to the two rules, so every object below is compared at one and the
same plan and one and the same signal law.

\emph{The weight leg.} Let $j$ be a Voice plan and $s$ a signal whose history is flagged under the
longer clock, $D_j(s;\tau,T) = 1$. By the clock equivalence, the filing lands by the horizon exactly
when the stake $T$ rounds before the horizon has reached the threshold, so
$B_j(s,H-T) \ge \tau$. Since $T' < T$ gives $H-T' > H-T$, and since the stake path of a Voice plan
is weakly increasing in the round, $B_j(s,H-T') \ge B_j(s,H-T) \ge \tau$. The engagement flag
$a_j = 1$ does not move with the clock. Reading the clock equivalence in the other direction at the
shorter clock gives $D_j(s;\tau,T') = 1$. Only Voice plans reach the threshold, so no other history
has to be checked, and the flagged cell at the longer clock is contained in the flagged cell at the
shorter clock. Taking probabilities under the common signal law, $\Omega(\tau,T') \ge
\Omega(\tau,T)$, so $1 - \Omega(\tau,T') \le 1 - \Omega(\tau,T)$. Both sides are strictly
positive because $\Omega < 1$ at each clock, so the quotient in \eqref{eq:clock-legs} is defined and
$0 < W_T \le 1$. This is clause~(i).

\emph{The ratio identity.} By Proposition~\ref{prop:factorisation} at both clocks, which the
statement assumes, the noise sensitivity of the expected engagement premium is the pooled share
times the pooled sensitivity. The factorisation hypotheses at the shorter clock, differentiability
of $\kappa \mapsto M_P(\kappa,\tau,T')$ included, make $\mathcal S_P(\kappa,\tau,T')$ exist, so
$C_T$ in \eqref{eq:clock-legs} is defined. At each clock,
\[
\mathcal S(\kappa,\tau,T) \;=\; \bigl(1-\Omega(\tau,T)\bigr)\,\mathcal S_P(\kappa,\tau,T),
\qquad
\mathcal S(\kappa,\tau,T') \;=\; \bigl(1-\Omega(\tau,T')\bigr)\,\mathcal S_P(\kappa,\tau,T') .
\]
The flagged weight does not move with $\kappa$, which is what licenses writing it as
$\Omega(\tau,\cdot)$ with no liquidity argument on either side. Since $1-\Omega(\tau,T) > 0$ and
$\mathcal S_P(\kappa,\tau,T) > 0$, the denominator $\mathcal S(\kappa,\tau,T)$ is strictly positive
and the quotient of the two displays is defined. Dividing and grouping the two factors gives
\eqref{eq:clock-ratio}, which is clause~(ii).

\emph{The equivalence.} Multiplying an inequality by the strictly positive number
$\mathcal S(\kappa,\tau,T)$ preserves it in both directions, so
$\mathcal S(\kappa,\tau,T') \le \mathcal S(\kappa,\tau,T)$ holds exactly when
$\mathcal S(\kappa,\tau,T')/\mathcal S(\kappa,\tau,T) \le 1$, and by \eqref{eq:clock-ratio} that
quotient is $W_T C_T$. This is \eqref{eq:clock-iff} and clause~(iii).
\end{proof}

\noindent
Because $W_T > 0$, the criterion \eqref{eq:clock-iff} can be read as a ceiling on the composition
leg: the shorter clock lowers noise sensitivity exactly when
\[
C_T \;\le\; \frac{1}{W_T} \;=\; \frac{1-\Omega(\tau,T)}{1-\Omega(\tau,T')} .
\]
By clause~(i) that ceiling is at least one, so $C_T \le 1$ is sufficient for attenuation at the
clock margin.
